\documentclass[12pt,a4paper]{article}
\usepackage[utf8]{inputenc}
\usepackage[french]{babel}
\usepackage{geometry}
\usepackage{graphicx}
\usepackage{xcolor}
\usepackage{booktabs}
\usepackage{longtable}
\usepackage{array}
\usepackage{enumitem}
\usepackage{fancyhdr}
\usepackage{titlesec}
\usepackage{hyperref}
\usepackage{amsmath}
\usepackage{amssymb}
\usepackage{tikz}
\usepackage{pgfplots}

\geometry{margin=2.5cm}
\pagestyle{fancy}
\fancyhf{}
\fancyhead[L]{\textbf{Tableau de Bord FDEM}}
\fancyhead[R]{\thepage}
\fancyfoot[C]{\textit{Indicateurs de Performance - Document Confidentiel}}

\definecolor{primaryblue}{RGB}{33, 150, 243}
\definecolor{secondaryblue}{RGB}{144, 202, 249}
\definecolor{darkblue}{RGB}{25, 118, 210}
\definecolor{lightgray}{RGB}{245, 245, 245}

\titleformat{\section}{\Large\bfseries\color{primaryblue}}{\thesection}{1em}{}
\titleformat{\subsection}{\large\bfseries\color{darkblue}}{\thesubsection}{1em}{}

\hypersetup{
    colorlinks=true,
    linkcolor=darkblue,
    filecolor=magenta,      
    urlcolor=cyan,
    pdftitle={Tableau de Bord FDEM - Indicateurs},
    pdfauthor={Direction Générale},
    pdfsubject={Indicateurs de Performance},
    pdfkeywords={Dashboard, KPI, Indicateurs, FDEM}
}

\begin{document}

\begin{titlepage}
    \centering
    \vspace*{2cm}
    
    {\Huge\bfseries\color{primaryblue} TABLEAU DE BORD FDEM}\\[0.5cm]
    {\Large\color{darkblue} Indicateurs de Performance Détaillés}\\[2cm]
    
    \begin{tikzpicture}
        \draw[primaryblue, line width=3pt] (0,0) -- (10,0);
        \draw[secondaryblue, line width=2pt] (0,-0.2) -- (10,-0.2);
    \end{tikzpicture}\\[2cm]
    
    {\large\textbf{Document de Référence pour la Direction Générale}}\\[1cm]
    {\large Analyse Complète des Indicateurs Existants et Proposés}\\[3cm]
    
    \begin{tabular}{ll}
        \textbf{Version:} & 1.0 \\
        \textbf{Date:} & \today \\
        \textbf{Statut:} & Confidentiel \\
        \textbf{Destinataire:} & Direction Générale
    \end{tabular}
    
    \vfill
    {\small\textit{Ce document contient des informations confidentielles destinées exclusivement à la Direction Générale}}
\end{titlepage}

\tableofcontents
\newpage

\section{Résumé Exécutif}

Ce document présente une analyse complète des indicateurs de performance (KPI) du système de gestion documentaire FDEM. Il comprend les indicateurs actuellement implémentés dans le tableau de bord ainsi que des propositions d'amélioration pour optimiser le pilotage stratégique.

\subsection{Objectifs du Document}
\begin{itemize}
    \item Inventorier les indicateurs existants
    \item Proposer de nouveaux indicateurs stratégiques
    \item Définir les métriques de performance clés
    \item Établir un référentiel pour le pilotage
\end{itemize}

\section{Indicateurs Existants - Analyse Détaillée}

\subsection{Indicateurs Principaux du Dashboard}

Le tableau de bord actuel présente plusieurs catégories d'indicateurs essentiels pour le suivi opérationnel :

\begin{longtable}{|p{4cm}|p{3cm}|p{7cm}|}
\hline
\rowcolor{lightgray}
\textbf{Indicateur} & \textbf{Type} & \textbf{Description} \\
\hline
\endhead

Total Documents & Quantitatif & Nombre total de documents dans le système \\
\hline
Documents Récents & Temporel & Documents ajoutés dans les dernières 24h/7 jours \\
\hline
Types de Produits & Catégoriel & Répartition par type de produit (Béton, Acier, etc.) \\
\hline
Structures Actives & Organisationnel & Nombre de structures organisationnelles \\
\hline
Projets en Cours & Opérationnel & Nombre de projets actuellement actifs \\
\hline
Utilisateurs Actifs & Ressources & Nombre d'utilisateurs connectés/actifs \\
\hline
Taux de Conformité & Qualité & Pourcentage de documents conformes aux standards \\
\hline
Documents Expirés & Risque & Documents nécessitant une mise à jour \\
\hline
\end{longtable}

\subsection{Indicateurs Hiérarchiques}

Le système propose une vue hiérarchique détaillée permettant l'analyse multi-niveaux :

\begin{enumerate}
    \item \textbf{Niveau 1 - Types de Produits}
    \begin{itemize}
        \item Béton Armé
        \item Structures Métalliques
        \item Matériaux Composites
        \item Équipements Spécialisés
    \end{itemize}
    
    \item \textbf{Niveau 2 - Produits Spécifiques}
    \begin{itemize}
        \item Sous-catégories par type
        \item Spécifications techniques
        \item Variantes et options
    \end{itemize}
    
    \item \textbf{Niveau 3 - Structures Organisationnelles}
    \begin{itemize}
        \item Départements
        \item Équipes projet
        \item Unités opérationnelles
    \end{itemize}
    
    \item \textbf{Niveau 4 - Sections et Subdivisions}
    \begin{itemize}
        \item Sections techniques
        \item Subdivisions N1, N2, N3
        \item Granularité maximale
    \end{itemize}
\end{enumerate}

\section{Indicateurs Proposés - Amélioration Stratégique}

\subsection{Indicateurs de Performance Opérationnelle}

\begin{longtable}{|p{4cm}|p{2cm}|p{2cm}|p{5cm}|}
\hline
\rowcolor{lightgray}
\textbf{Indicateur} & \textbf{Priorité} & \textbf{Fréquence} & \textbf{Objectif Stratégique} \\
\hline
\endhead

Temps Moyen de Traitement & Haute & Quotidien & Optimiser l'efficacité opérationnelle \\
\hline
Taux de Résolution Premier Contact & Haute & Hebdomadaire & Améliorer la satisfaction client \\
\hline
Charge de Travail par Équipe & Moyenne & Quotidien & Équilibrer les ressources \\
\hline
Délai Moyen de Validation & Haute & Hebdomadaire & Accélérer les processus \\
\hline
Taux d'Utilisation Système & Moyenne & Mensuel & Optimiser l'adoption technologique \\
\hline
\end{longtable}

\subsection{Indicateurs de Qualité et Conformité}

\begin{longtable}{|p{4cm}|p{2cm}|p{2cm}|p{5cm}|}
\hline
\rowcolor{lightgray}
\textbf{Indicateur} & \textbf{Priorité} & \textbf{Fréquence} & \textbf{Objectif Stratégique} \\
\hline
\endhead

Taux de Non-Conformité & Critique & Quotidien & Maintenir les standards qualité \\
\hline
Nombre d'Audits Réalisés & Haute & Mensuel & Assurer la conformité réglementaire \\
\hline
Délai de Correction Anomalies & Haute & Hebdomadaire & Réduire les risques opérationnels \\
\hline
Taux de Satisfaction Qualité & Moyenne & Trimestriel & Améliorer la perception qualité \\
\hline
Coût de la Non-Qualité & Critique & Mensuel & Optimiser les coûts qualité \\
\hline
\end{longtable}

\subsection{Indicateurs Financiers et Budgétaires}

\begin{longtable}{|p{4cm}|p{2cm}|p{2cm}|p{5cm}|}
\hline
\rowcolor{lightgray}
\textbf{Indicateur} & \textbf{Priorité} & \textbf{Fréquence} & \textbf{Objectif Stratégique} \\
\hline
\endhead

Coût par Document Traité & Haute & Mensuel & Optimiser l'efficacité économique \\
\hline
ROI des Projets & Critique & Trimestriel & Maximiser le retour sur investissement \\
\hline
Budget Consommé vs Prévu & Critique & Mensuel & Contrôler les dépenses \\
\hline
Coût de Maintenance Système & Moyenne & Mensuel & Optimiser les coûts IT \\
\hline
Économies Réalisées & Haute & Trimestriel & Mesurer les gains d'efficacité \\
\hline
\end{longtable}

\section{Indicateurs de Ressources Humaines}

\subsection{Performance et Productivité}

\begin{itemize}
    \item \textbf{Productivité par Employé} : Documents traités par personne/jour
    \item \textbf{Taux d'Absentéisme} : Impact sur la performance globale
    \item \textbf{Temps de Formation} : Investissement en développement des compétences
    \item \textbf{Taux de Rotation} : Stabilité des équipes
    \item \textbf{Satisfaction Employés} : Enquêtes internes régulières
\end{itemize}

\subsection{Compétences et Développement}

\begin{itemize}
    \item \textbf{Niveau de Certification} : Pourcentage d'employés certifiés
    \item \textbf{Heures de Formation} : Investissement formation par employé
    \item \textbf{Évaluations Performance} : Scores moyens d'évaluation
    \item \textbf{Promotions Internes} : Taux de mobilité ascendante
\end{itemize}

\section{Indicateurs Technologiques et Innovation}

\subsection{Performance Système}

\begin{longtable}{|p{4cm}|p{3cm}|p{6cm}|}
\hline
\rowcolor{lightgray}
\textbf{Métrique} & \textbf{Seuil Critique} & \textbf{Action Corrective} \\
\hline
\endhead

Temps de Réponse Système & > 3 secondes & Optimisation infrastructure \\
\hline
Disponibilité Système & < 99.5\% & Renforcement redondance \\
\hline
Taux d'Erreur & > 1\% & Investigation technique urgente \\
\hline
Capacité Stockage & > 85\% & Extension infrastructure \\
\hline
Sécurité - Tentatives Intrusion & > 0 & Audit sécurité immédiat \\
\hline
\end{longtable}

\subsection{Innovation et Évolution}

\begin{itemize}
    \item \textbf{Nouvelles Fonctionnalités Déployées} : Rythme d'innovation
    \item \textbf{Adoption Nouvelles Technologies} : Taux d'utilisation
    \item \textbf{Suggestions Amélioration} : Nombre d'idées proposées
    \item \textbf{Projets R\&D} : Investissement en recherche
\end{itemize}

\section{Indicateurs de Risque et Sécurité}

\subsection{Gestion des Risques Opérationnels}

\begin{enumerate}
    \item \textbf{Risques Identifiés}
    \begin{itemize}
        \item Nombre de risques par catégorie
        \item Niveau de criticité moyen
        \item Délai de résolution
    \end{itemize}
    
    \item \textbf{Incidents de Sécurité}
    \begin{itemize}
        \item Nombre d'incidents par mois
        \item Temps de résolution moyen
        \item Impact sur les opérations
    \end{itemize}
    
    \item \textbf{Conformité Réglementaire}
    \begin{itemize}
        \item Taux de conformité aux normes
        \item Nombre d'audits réussis
        \item Actions correctives en cours
    \end{itemize}
\end{enumerate}

\section{Indicateurs Client et Satisfaction}

\subsection{Satisfaction Client}

\begin{longtable}{|p{4cm}|p{2cm}|p{2cm}|p{4cm}|}
\hline
\rowcolor{lightgray}
\textbf{Indicateur} & \textbf{Cible} & \textbf{Actuel} & \textbf{Tendance} \\
\hline
\endhead

Score Satisfaction Globale & 4.5/5 & 4.2/5 & \textcolor{green}{↗} \\
\hline
Taux de Recommandation (NPS) & > 70 & 65 & \textcolor{orange}{→} \\
\hline
Délai Réponse Réclamations & < 24h & 18h & \textcolor{green}{↗} \\
\hline
Taux de Résolution Premier Contact & > 80\% & 75\% & \textcolor{red}{↘} \\
\hline
\end{longtable}

\subsection{Engagement Client}

\begin{itemize}
    \item \textbf{Fréquence d'Utilisation} : Connexions par utilisateur/mois
    \item \textbf{Durée Sessions} : Temps moyen passé dans l'application
    \item \textbf{Fonctionnalités Utilisées} : Taux d'adoption par fonctionnalité
    \item \textbf{Feedback Utilisateurs} : Nombre de retours constructifs
\end{itemize}

\section{Tableaux de Bord Recommandés}

\subsection{Dashboard Exécutif (Direction Générale)}

\begin{enumerate}
    \item \textbf{Vue d'Ensemble Stratégique}
    \begin{itemize}
        \item KPI financiers principaux
        \item Indicateurs de performance globale
        \item Alertes critiques
        \item Tendances à long terme
    \end{itemize}
    
    \item \textbf{Analyse Comparative}
    \begin{itemize}
        \item Performance vs objectifs
        \item Benchmarking sectoriel
        \item Évolution historique
        \item Projections futures
    \end{itemize}
\end{enumerate}

\subsection{Dashboard Opérationnel (Managers)}

\begin{enumerate}
    \item \textbf{Suivi Quotidien}
    \begin{itemize}
        \item Charge de travail équipes
        \item Indicateurs de productivité
        \item Alertes opérationnelles
        \item Planning et ressources
    \end{itemize}
    
    \item \textbf{Analyse Performance}
    \begin{itemize}
        \item Métriques par équipe
        \item Comparaisons inter-services
        \item Identification des goulots d'étranglement
        \item Plans d'action
    \end{itemize}
\end{enumerate}

\section{Recommandations d'Implémentation}

\subsection{Phase 1 : Indicateurs Critiques (0-3 mois)}

\begin{enumerate}
    \item Mise en place des indicateurs financiers de base
    \item Implémentation du suivi temps réel
    \item Création des alertes automatiques
    \item Formation des équipes utilisatrices
\end{enumerate}

\subsection{Phase 2 : Indicateurs Avancés (3-6 mois)}

\begin{enumerate}
    \item Déploiement des métriques de qualité
    \item Intégration des indicateurs RH
    \item Mise en place du reporting automatisé
    \item Optimisation des processus
\end{enumerate}

\subsection{Phase 3 : Innovation et Prédictif (6-12 mois)}

\begin{enumerate}
    \item Implémentation d'analytics avancés
    \item Développement d'indicateurs prédictifs
    \item Intelligence artificielle pour l'analyse
    \item Tableaux de bord personnalisés
\end{enumerate}

\section{Métriques de Succès}

\subsection{Critères d'Évaluation}

\begin{longtable}{|p{4cm}|p{3cm}|p{6cm}|}
\hline
\rowcolor{lightgray}
\textbf{Objectif} & \textbf{Métrique} & \textbf{Cible 2024} \\
\hline
\endhead

Amélioration Efficacité & Temps de traitement & -25\% vs 2023 \\
\hline
Réduction Coûts & Coût par document & -15\% vs 2023 \\
\hline
Satisfaction Client & Score NPS & > 75 \\
\hline
Qualité Processus & Taux de non-conformité & < 2\% \\
\hline
Innovation & Nouvelles fonctionnalités & 12 par an \\
\hline
\end{longtable}

\section{Conclusion et Prochaines Étapes}

\subsection{Synthèse}

Le système actuel dispose d'une base solide d'indicateurs opérationnels. Les améliorations proposées permettront :

\begin{itemize}
    \item Une vision stratégique renforcée
    \item Un pilotage plus précis des performances
    \item Une anticipation des risques améliorée
    \item Une prise de décision basée sur des données fiables
\end{itemize}

\subsection{Plan d'Action Immédiat}

\begin{enumerate}
    \item \textbf{Validation Direction} : Approbation des indicateurs prioritaires
    \item \textbf{Ressources Techniques} : Allocation équipe développement
    \item \textbf{Formation Utilisateurs} : Programme de formation adapté
    \item \textbf{Calendrier Déploiement} : Planning détaillé par phase
\end{enumerate}

\subsection{Bénéfices Attendus}

\begin{itemize}
    \item \textbf{Amélioration de 30\%} de l'efficacité opérationnelle
    \item \textbf{Réduction de 20\%} des coûts de non-qualité
    \item \textbf{Augmentation de 25\%} de la satisfaction client
    \item \textbf{ROI estimé} : 250\% sur 24 mois
\end{itemize}

\vfill

\begin{center}
\textit{Document préparé pour la Direction Générale FDEM}\\
\textit{Confidentiel - Ne pas diffuser}
\end{center}

\end{document}